\chapter{Pecahan dan Pangkat}\label{ch:g9:fractions}

\omterm{def:g9:fractions:fraction}{Pecahan} dan \omterm{def:g9:fractions:power}{pangkat} adalah alat perhitungan sehari-hari: membagi
besaran, membandingkan proporsi, menulis bilangan yang sangat besar
atau sangat kecil. Bab ini memantapkan aturan berhitung dengannya ---
dengan setiap langkah perantaranya ditulis lengkap --- lalu
memperkenalkan \omterm{def:g9:fractions:scientific}{notasi ilmiah}.

\section{Berhitung dengan pecahan}

\begin{definition}[Pecahan]\label{def:g9:fractions:fraction}
\emph{Pecahan}\index{pecahan} $\dfrac ab$ (dengan $b \neq 0$) adalah
\omterm{def:g3:division:remainder}{hasil bagi} $a$ oleh $b$. Dua pecahan sama ketika yang satu diperoleh
dari yang lain dengan mengalikan (atau membagi) \omterm{def:g4:fractions:def}{pembilang} dan
\omterm{def:g4:fractions:def}{penyebutnya} dengan bilangan bukan \omterm{def:g1:counting:zero}{nol} yang sama:
\[
\frac ab = \frac{a \times k}{b \times k} \qquad (k \neq 0).
\]
\end{definition}

\begin{example}\label{ex:g9:fractions:simplify}
Sederhanakan $\dfrac{42}{56}$ selangkah demi selangkah:
\[
\frac{42}{56} = \frac{42 \div 2}{56 \div 2} = \frac{21}{28}
= \frac{21 \div 7}{28 \div 7} = \frac34 .
\]
(Pada \cref{ch:g9:arith} faktor persekutuan terbesar akan
mengerjakannya dalam satu langkah.)
\end{example}

\begin{method}[Menjumlahkan dan mengurangkan pecahan]\label{met:g9:fractions:add}
\begin{enumerate}
  \item Taruhlah kedua \omterm{def:g9:fractions:fraction}{pecahannya} pada \emph{\omterm{def:g4:fractions:def}{penyebut} yang sama}
        (yaitu \omterm{def:g5:division:multiple}{kelipatan} bersama kedua \omterm{def:g4:fractions:def}{penyebutnya});
  \item jumlahkan atau kurangkan \omterm{def:g4:fractions:def}{pembilangnya}, dengan menahan
        \omterm{def:g4:fractions:def}{penyebutnya};
  \item sederhanakan hasilnya kalau mungkin.
\end{enumerate}
\end{method}

\begin{example}\label{ex:g9:fractions:add}
Hitunglah $\dfrac56 + \dfrac34$. \omterm{def:g4:fractions:def}{Penyebut} bersama dari $6$ dan $4$
adalah $12$:
\[
\frac56 + \frac34
= \frac{5 \times 2}{6 \times 2} + \frac{3 \times 3}{4 \times 3}
= \frac{10}{12} + \frac{9}{12}
= \frac{19}{12}.
\]
Hitunglah $2 - \dfrac37$. Tulislah $2$ sebagai \omterm{def:g9:fractions:fraction}{pecahan} berpenyebut $7$:
\[
2 - \frac37 = \frac{14}{7} - \frac37 = \frac{11}{7}.
\]
\end{example}

\begin{omfigure}
\begin{tikzpicture}[scale=0.55]
  \foreach \i in {0,...,11}
    \draw[omExo] ({mod(\i,12)},0) rectangle ({mod(\i,12)+1},1);
  \foreach \i in {0,...,9}
    \fill[omDef!35] (\i,0) rectangle (\i+1,1);
  \draw[thick] (0,0) rectangle (12,1);
  \node[below, font=\small] at (6,-0.1)
    {$\frac56 = \frac{10}{12}$: 10 kotak dari 12};
  \begin{scope}[yshift=-2.4cm]
  \foreach \i in {0,...,11}
    \draw[omExo] (\i,0) rectangle (\i+1,1);
  \foreach \i in {0,...,8}
    \fill[omThm!35] (\i,0) rectangle (\i+1,1);
  \draw[thick] (0,0) rectangle (12,1);
  \node[below, font=\small] at (6,-0.1)
    {$\frac34 = \frac{9}{12}$: 9 kotak dari 12};
  \end{scope}
\end{tikzpicture}

{\small Mengapa perlu \omterm{def:g4:fractions:def}{penyebut} yang sama: memotong kedua keseluruhannya
menjadi $12$ kotak yang sama membuat \omterm{def:g9:fractions:fraction}{pecahannya} dapat dibandingkan dan
dijumlahkan --- bersama-sama, $10 + 9 = 19$ perdua belas.}
\end{omfigure}

\begin{method}[Mengalikan dan membagi pecahan]\label{met:g9:fractions:mult}
\begin{enumerate}
  \item Untuk mengalikan, kalikan \omterm{def:g4:fractions:def}{pembilangnya} bersama dan
        \omterm{def:g4:fractions:def}{penyebutnya} bersama: $\dfrac ab \times \dfrac cd = \dfrac{ac}{bd}$
        (sederhanakan \emph{sebelum} mengalikan bilamana mungkin);
  \item untuk membagi, kalikan dengan \omterm{def:g8:fractions:inverse}{kebalikan} \omterm{def:g5:division:multiple}{pembaginya}:
        $\dfrac ab \div \dfrac cd = \dfrac ab \times \dfrac dc$.
\end{enumerate}
\end{method}

\begin{example}\label{ex:g9:fractions:mult}
\[
\frac{7}{15} \times \frac{25}{14}
= \frac{7 \times 25}{15 \times 14}
= \frac{7 \times 25}{15 \times 2 \times 7}
= \frac{25}{30} = \frac56,
\]
yang kita sederhanakan dengan faktor bersamanya $7$, lalu dengan $5$.
Dan sebuah pembagian:
\[
\frac{3}{4} \div \frac{9}{8}
= \frac34 \times \frac89
= \frac{3 \times 8}{4 \times 9}
= \frac{24}{36} = \frac23 .
\]
\end{example}

\section{Pangkat}

\begin{definition}[Pangkat]\label{def:g9:fractions:power}
Untuk bilangan real $a$ dan bilangan bulat positif $n$,
\emph{pangkat}\index{pangkat} $a^n$ adalah \omterm{def:g2:mult:def}{hasil kali} $n$ faktor yang
masing-masing sama dengan $a$:
\[
a^n = \underbrace{a \times a \times \dots \times a}_{n \text{ faktor}} .
\]
Menurut kesepakatan $a^0 = 1$ (untuk $a \neq 0$), dan \omterm{def:g8:powers:def}{eksponen} negatif
menyatakan \omterm{def:g8:fractions:inverse}{kebalikannya}:
inverses:
\[
a^{-n} = \frac{1}{a^n}.
\]
\end{definition}

\begin{example}
$2^4 = 2 \times 2 \times 2 \times 2 = 16$; \quad
$(-3)^2 = 9$ tetapi $-3^2 = -(3^2) = -9$ (\omterm{def:g8:powers:def}{eksponennya} mengikat lebih
dahulu daripada tanda minusnya); \quad
$10^{-3} = \frac{1}{1000} = 0.001$; \quad
$5^1 = 5$ dan $5^0 = 1$.
\end{example}

\begin{theorem}[Aturan eksponen]\label{thm:g9:fractions:rules}
Untuk semua $a$, $b$ yang bukan \omterm{def:g1:counting:zero}{nol} dan semua bilangan bulat $m$, $n$:
\[
a^m \times a^n = a^{m+n},
\qquad
\frac{a^m}{a^n} = a^{m-n},
\qquad
\left(a^m\right)^n = a^{mn},
\qquad
(ab)^n = a^n b^n .
\]
\end{theorem}

\begin{proof}
Untuk \omterm{def:g8:powers:def}{eksponen} positif, bilanglah faktornya. Untuk aturan pertama:
$a^m \times a^n$ adalah $m$ faktor $a$ yang disusul $n$ faktor $a$,
seluruhnya $m + n$ faktor. Untuk yang ketiga:
$\left(a^m\right)^n$ mengulang $n$ kali satu blok berisi $m$ faktor,
sehingga memberi $mn$ faktor. Kedua aturan lainnya serupa; perluasan
ke \omterm{def:g8:powers:def}{eksponen} \omterm{def:g1:counting:zero}{nol} dan negatif diperiksa dari $a^{-n} = \frac{1}{a^n}$
(dan itu membuat kesepakatan $a^0 = 1$ menjadi satu-satunya pilihan
yang taat asas, karena $a^n \times a^0$ harus sama dengan
$a^{n+0} = a^n$).
\end{proof}

\begin{example}\label{ex:g9:fractions:rules}
Sederhanakan selangkah demi selangkah:
\[
\frac{3^7 \times 3^{-2}}{3^4}
= \frac{3^{7 + (-2)}}{3^4}
= \frac{3^5}{3^4}
= 3^{5-4} = 3^1 = 3 .
\]
Dan dengan dua huruf: $(2a^3)^2 \times a = 4 a^6 \times a = 4a^7$.
\end{example}

\section{Pangkat sepuluh dan notasi ilmiah}

\begin{proposition}[Pangkat sepuluh]\label{prop:g9:fractions:ten}
Untuk bilangan bulat positif $n$: $10^n$ ditulis ``$1$ yang disusul
$n$ \omterm{def:g1:counting:zero}{nol}'', dan $10^{-n} = 0.0\dots01$ dengan $1$-nya pada tempat
desimal yang ke-$n$. Mengalikan \omterm{def:g6:decimals:places}{bilangan desimal} dengan $10^n$
(berturut-turut $10^{-n}$) menggeser \omterm{def:g4:decimals:point}{titik desimalnya} $n$ tempat ke
kanan (berturut-turut ke kiri).
\end{proposition}

\begin{definition}[Notasi ilmiah]\label{def:g9:fractions:scientific}
\emph{Notasi ilmiah}\index{notasi ilmiah} sebuah \omterm{def:g6:decimals:places}{bilangan desimal} yang
positif adalah satu-satunya cara menuliskannya sebagai
\[
a \times 10^n,
\qquad\text{dengan } 1 \leq a < 10 \text{ dan } n \text{ bilangan bulat.}
\]
\end{definition}

\begin{example}\label{ex:g9:fractions:scientific}
Jarak Bumi--Matahari kira-kira $149\,600\,000$ km $= 1.496 \times
10^8$ km. Ukuran sebuah molekul air kira-kira
$0.000\,000\,000\,28$ m $= 2.8 \times 10^{-10}$ m. Untuk berhitung
dengan bilangan seperti itu, kelompokkan bagian desimalnya dan \omterm{def:g9:fractions:power}{pangkat}
sepuluhnya:
\[
(3 \times 10^5) \times (2.5 \times 10^{-2})
= (3 \times 2.5) \times 10^{5 + (-2)}
= 7.5 \times 10^{3}.
\]
\end{example}

\begin{method}[Menaruh bilangan dalam notasi ilmiah]\label{met:g9:fractions:scientific}
\begin{enumerate}
  \item Geserlah \omterm{def:g4:decimals:point}{titik desimalnya} sampai tepat sesudah angka
        pertama yang bukan \omterm{def:g1:counting:zero}{nol}; ini memberi faktor $a$ dengan
        $1 \leq a < 10$;
  \item bilanglah berapa tempat titiknya bergeser: bilangan itulah
        $n$, positif kalau bilangan semulanya $\geq 10$, negatif
        kalau bilangan semulanya $< 1$;
  \item periksa: $a \times 10^n$ harus menghasilkan bilangan semulanya.
\end{enumerate}
\end{method}

\section{Latihan}

\begin{exercise}[$\star$]\label{exo:g9:fractions:1}
Hitunglah lalu berikan hasilnya sebagai \omterm{def:g9:fractions:fraction}{pecahan} yang paling sederhana:
\[
\frac13 + \frac15, \qquad
\frac72 - \frac53, \qquad
\frac{5}{12} + \frac38, \qquad
3 - \frac45 .
\]
\end{exercise}

\begin{exercise}[$\star$]\label{exo:g9:fractions:2}
Hitunglah lalu sederhanakan:
\[
\frac59 \times \frac{27}{35}, \qquad
\frac{8}{15} \div \frac{4}{5}, \qquad
\frac23 \times \frac34 \times \frac45 .
\]
\end{exercise}

\begin{exercise}[$\star$]\label{exo:g9:fractions:3}
Hitunglah:
\[
2^5, \qquad (-2)^4, \qquad -2^4, \qquad 4^{-2}, \qquad
\left(\tfrac23\right)^3, \qquad 7^0 .
\]
\end{exercise}

\begin{exercise}[$\star$]\label{exo:g9:fractions:4}
Tulislah sebagai satu \omterm{def:g9:fractions:power}{pangkat}:
\[
5^3 \times 5^6, \qquad
\frac{7^9}{7^5}, \qquad
\left(2^3\right)^4, \qquad
\frac{3^2 \times 3^7}{3^5}, \qquad
2^5 \times 4 .
\]
\end{exercise}

\begin{exercise}[$\star$]\label{exo:g9:fractions:5}
Tulislah dalam \omterm{def:g9:fractions:scientific}{notasi ilmiah}: $45\,000$; $0.0072$; $302.5$;
$0.000\,000\,9$; dua belas juta.
\end{exercise}

\begin{exercise}[$\star\star$]\label{exo:g9:fractions:6}
Hitunglah, dengan memberikan hasilnya dalam \omterm{def:g9:fractions:scientific}{notasi ilmiah}:
\[
(2 \times 10^7) \times (4.5 \times 10^{-3}),
\qquad
\frac{9 \times 10^5}{3 \times 10^{-2}},
\qquad
(5 \times 10^3)^2 .
\]
\end{exercise}

\begin{exercise}[$\star\star$]\label{exo:g9:fractions:7}
Sebuah tangki terisi $\frac35$. Sesudah ditambah $30$ liter tangkinya
terisi $\frac34$. Berapa \omterm{def:g2:measure:units}{daya tampung} tangkinya? (Nyatakan lebih
dahulu bagian tangki yang ditambahkan.)
\end{exercise}

\begin{exercise}[$\star\star$]\label{exo:g9:fractions:8}
Alia memakan $\frac14$ sebuah kue, lalu Beni memakan $\frac13$ \omterm{def:g3:division:remainder}{sisanya},
lalu Cinta memakan $\frac12$ dari yang masih tersisa. Berapa bagian
kuenya yang tertinggal pada akhirnya? (Hitunglah \omterm{def:g3:division:remainder}{sisanya} sesudah tiap
langkahnya.)
\end{exercise}

\begin{exercise}[$\star\star$]\label{exo:g9:fractions:9}
Cahaya menempuh $3 \times 10^8$ meter tiap detik. Matahari berjarak
kira-kira $1.5 \times 10^{11}$ m. Berapa lama sinar matahari sampai
kepada kita? Berikan jawabannya dalam detik (tanpa perlu \omterm{def:g9:fractions:scientific}{notasi
ilmiah}), lalu dalam menit dan detik.
\end{exercise}

\begin{exercise}[$\star\star\star$]\label{exo:g9:fractions:10}
Mana yang lebih besar, $2^{30}$ atau $3^{20}$? (Petunjuk: tulislah
keduanya sebagai \omterm{def:g9:fractions:power}{pangkat} dengan \omterm{def:g8:powers:def}{eksponen} $10$ memakai
$\left(a^m\right)^n = a^{mn}$, lalu bandingkan $2^3$ dengan
$3^2$.)
\end{exercise}

\section{Soal: Sandi rahasia desimal berulang}

\begin{problem}[{Soal akhir pekan --- setiap pecahan punya penulisan
desimal yang berhenti atau berulang, setiap desimal berulang adalah
pecahan, dan $0.999\ldots$ sama persis dengan $1$}]
\label{pb:g9:fractions:1}
Bagilah $3$ dengan $11$ dan angkanya jatuh menjadi nyanyian:
$0.272727\ldots$ selamanya. Bagilah $1$ dengan $4$ dan angkanya
berhenti mati: $0.25$. Soal ini membuktikan bahwa hanya kedua
perilaku itulah yang mungkin bagi sebuah \omterm{def:g9:fractions:fraction}{pecahan} --- lalu membalikkan
mesinnya: diberikan desimal berulang mana pun, mesinnya membangun
kembali \omterm{def:g9:fractions:fraction}{pecahan} yang bersembunyi di baliknya. Di sepanjang jalannya
soal ini menyelesaikan, sekali untuk selamanya, kesamaan yang paling
diperdebatkan dalam matematika sekolah.

\medskip
\noindent\textbf{Bagian I --- Dari \omterm{def:g9:fractions:fraction}{pecahan} ke desimal.}
\begin{enumerate}
  \item Hitunglah dengan pembagian bersusun penulisan desimal
        $\frac14$, $\frac13$, $\frac56$ dan $\frac{3}{11}$.
        Mana yang berhenti, mana yang berulang, dan dengan blok
        berulang yang mana?
  \item Jelaskan mengapa penulisan desimal \omterm{def:g9:fractions:fraction}{pecahan} \emph{apa pun}
        $\frac ab$ pasti berhenti atau akhirnya berulang. (Pada
        pembagian bersusun dengan $b$, nilai apa saja yang dapat
        diambil \omterm{def:g3:division:remainder}{sisa} yang berturut-turut? Apa yang terjadi begitu
        sebuah \omterm{def:g3:division:remainder}{sisa} muncul untuk kedua kalinya?)
  \item Doronglah pembagian $1$ oleh $7$ cukup jauh untuk menemukan
        blok berulangnya. Berapa panjangnya, dan bagaimana
        perbandingan panjang itu dengan batas yang dijanjikan
        pertanyaan 2?
  \item Beberapa \omterm{def:g9:fractions:fraction}{pecahan} berhenti karena \omterm{def:g4:fractions:def}{penyebutnya} dapat
        digelembungkan menjadi \omterm{def:g9:fractions:power}{pangkat} sepuluh: $\frac14 =
        \frac{25}{100}$, $\frac{3}{8} = \frac{375}{1000}$.
        Jelaskan kriterianya, dan mengapa tidak ada pengali
        bilangan bulat yang dapat mengubah $3$, $6$ atau $11$
        menjadi $10$, $100$,
        $1000, \dots$ (angka apa yang akan mengakhiri \omterm{def:g2:mult:def}{hasil kalinya}?).
  \item Tanpa membagi, pilahlah $\frac{9}{40}$, $\frac{5}{12}$ dan
        $\frac{33}{50}$ menjadi ``berhenti'' dan ``berulang'', lalu
        berikan penulisan desimal kedua \omterm{def:g9:fractions:fraction}{pecahan} yang berhenti.
\end{enumerate}

\medskip
\noindent\textbf{Bagian II --- Dari desimal berulang kembali ke
\omterm{def:g9:fractions:fraction}{pecahannya}.}
\begin{enumerate}[resume]
  \item Trik penggeseran. Misalkan $x = 0.7777\ldots$ Hitunglah $10x$,
        lalu $10x - x$, lalu simpulkan $x$ sebagai \omterm{def:g9:fractions:fraction}{pecahan}.
        Periksalah dengan membagi.
  \item Misalkan $x = 0.727272\ldots$ \omterm{def:g9:fractions:power}{Pangkat} sepuluh yang mana yang
        menggeser $x$ tepat satu blok berulang? Pakailah itu untuk
        menulis $x$ sebagai \omterm{def:g9:fractions:fraction}{pecahan} yang paling sederhana.
  \item Sebuah keadaan campuran: $x = 0.58333\ldots$ (yang berulang
        hanya $3$-nya). Hitunglah $1000x - 100x$ lalu simpulkan $x$
        sebagai \omterm{def:g9:fractions:fraction}{pecahan} yang paling sederhana.
  \item Yang terkenal: terapkan trik penggeserannya pada
        $x = 0.9999\ldots$ \omterm{def:g9:fractions:fraction}{Pecahan} apa --- \emph{bilangan} apa ---
        yang kamu temukan? Rukunkan hasilnya dengan
        gerak hatimu, sambil mengingat tetangga hantu pada
        \cref{pb:g6:decimals:1} dan \omterm{def:g3:division:remainder}{sisa} yang mengerut pada
        \cref{pb:g6:fractions:1}: apakah $0.999\ldots$ ``tepat di
        bawah'' $1$, atau nama lain untuk $1$?
  \item Ubahlah $2.454545\ldots$ dan $0.123123123\ldots$ menjadi
        \omterm{def:g9:fractions:fraction}{pecahan} yang paling sederhana.
\end{enumerate}

\medskip
\noindent\textbf{Bagian III --- Kamus yang dilengkapi.}
\begin{enumerate}[resume]
  \item Pertanyaan 6--10 menyarankan sebuah kamus: blok berisi $k$
        angka yang berulang sejak \omterm{def:g4:decimals:point}{titik desimalnya} sama dengan blok
        itu dibagi $k$ buah sembilan ($0.727272\ldots = \frac{72}{99}$).
        Pakailah kamusnya untuk meramalkan $0.037037\ldots$ sebagai
        \omterm{def:g9:fractions:fraction}{pecahan}, sederhanakan --- lalu takjublah: \omterm{def:g9:fractions:fraction}{pecahan} yang
        sederhananya mengejutkan mana yang menyanyikan ``$037$''?
        Periksalah dengan pembagian.
  \item Gabungkan kamusnya dengan sebuah penggeseran: tulislah
        $0.0272727\ldots$ sebagai \omterm{def:g9:fractions:fraction}{pecahan} yang paling sederhana.
  \item Bilangan $0.101001000100001\ldots$ (satu $1$, lalu satu
        $0$, lalu satu $1$, lalu dua $0$, lalu satu $1$, lalu tiga
        $0$, dan seterusnya) tidak pernah berhenti. Apakah ia
        desimal yang berulang? Lalu apa kata Bagian I tentangnya
        --- dapatkah ia menjadi \omterm{def:g9:fractions:fraction}{pecahan}? (Kamu baru saja bertemu
        bilangan \emph{irasional} pertamamu yang terbukti; yang
        paling terkenal ditangkap pada \cref{ch:g9:sqrt}.)
  \item Kira-kira berapa angka yang dipunyai $2^{30}$? Pakailah
        $2^{10} = 1\,024 \approx 1.024 \times 10^3$ dan aturan
        \omterm{def:g8:powers:def}{eksponennya} (\cref{thm:g9:fractions:rules}) untuk menulis
        $2^{30}$ dalam \omterm{def:g9:fractions:scientific}{notasi ilmiah} (secara hampiran), lalu
        simpulkan. (Bandingkan \cref{exo:g9:fractions:10}.)
  \item Penutupnya: tulislah $0.20262026\ldots$ sebagai \omterm{def:g9:fractions:fraction}{pecahan},
        lalu nyatakan pesan moral dua arah seluruh soalnya:
        penulisan desimal yang mana yang berupa \omterm{def:g9:fractions:fraction}{pecahan}, dan
        \omterm{def:g9:fractions:fraction}{pecahan} yang mana yang punya penulisan desimal yang mana?

\end{enumerate}
\end{problem}
